% ==========================================================
% GENERATED FILE --- DO NOT EDIT.
% Source of truth: 01_law_of_graphic_equalisation_theory_referenced.tex
% Statement-only content: 01s.spec
% Regenerate, from the repository root:
% python3 tools/build_statement.py papers/formal/01s.spec
% ==========================================================
\documentclass[11pt,a4paper]{article}
\usepackage[utf8]{inputenc}
\usepackage[T1]{fontenc}
\usepackage{lmodern,amsmath,amssymb,amsthm,mathtools,array,booktabs,enumitem,microtype}
\usepackage{natbib}
\usepackage[margin=1in]{geometry}
\usepackage{fancyhdr}
\usepackage{xr-hyper}
\usepackage{hyperref}
\externaldocument[IS-]{00_is}
\externaldocument[PH-]{01a_graphic_equalisation_physics_grounding}
\externaldocument[MS-]{01b_graphic_equalisation_mind_and_self_grounding}
\externaldocument[SM-]{02_seldon_manifolds_theory_referenced}
\externaldocument[SG-]{02a_seldon_manifolds_psychological_grounding}
\externaldocument[NM-]{03_law_of_narrative_momentum_theory_referenced}
\externaldocument[NV-]{04_narrativium_theory_referenced}
\externaldocument[FR-]{05_frames_relation_and_the_empty_limit}
\externaldocument[AS-]{06_assessing_an_account_of_this_kind}
\externaldocument[AD-]{09_adaptability}
\newtheorem{axiom}{Axiom}
\newtheorem{definition}{Definition}[section]
\newtheorem{theorem}[definition]{Theorem}
\newtheorem{proposition}[definition]{Proposition}
\newtheorem{schema}[definition]{Schema}
\newtheorem{hypothesis}[definition]{Hypothesis}
\newtheorem{corollary}[definition]{Corollary}
\theoremstyle{remark}\newtheorem{remark}[definition]{Remark}
\newcommand{\Poss}{\mathcal{P}}
\newcommand{\Coll}{\mathcal{C}}
\newcommand{\Hist}{\mathcal{H}}
\pagestyle{fancy}
\fancyhf{}
\fancyfoot[C]{\thepage}
\fancyfoot[R]{\scriptsize\PaperFooterStamp}
\renewcommand{\headrulewidth}{0pt}
\fancypagestyle{plain}{\fancyhf{}\fancyfoot[C]{\thepage}%
  \fancyfoot[R]{\scriptsize\PaperFooterStamp}\renewcommand{\headrulewidth}{0pt}}
% GENERATED by tools/build_paper_metadata.py from paper_metadata.json.
% DO NOT EDIT: edit paper_metadata.json, then regenerate this file.
\providecommand{\PaperDraft}{draft-2026-09-06.1}
\providecommand{\PaperDisplayDate}{6 September 2026}
\expandafter\def\csname PaperRevision@00_is\endcsname{129}
\expandafter\def\csname PaperRevision@01_law_of_graphic_equalisation_theory_referenced\endcsname{294}
\expandafter\def\csname PaperRevision@01a_graphic_equalisation_physics_grounding\endcsname{47}
\expandafter\def\csname PaperRevision@01b_graphic_equalisation_mind_and_self_grounding\endcsname{47}
\expandafter\def\csname PaperRevision@01c_graphic_equalisation_cosmology_grounding\endcsname{11}
\expandafter\def\csname PaperRevision@01d_graphic_equalisation_born_recovery\endcsname{6}
\expandafter\def\csname PaperRevision@01e_graphic_equalisation_communication_grounding\endcsname{6}
\expandafter\def\csname PaperRevision@01f_graphic_equalisation_qualia_and_consciousness_grounding\endcsname{3}
\expandafter\def\csname PaperRevision@01g_graphic_equalisation_total_perspective_vortex_grounding\endcsname{4}
\expandafter\def\csname PaperRevision@01h_graphic_equalisation_information_gravity_quantum_reconciliation\endcsname{5}
\expandafter\def\csname PaperRevision@01i_graphic_equalisation_passive_continuation_momentum_gravity\endcsname{12}
\expandafter\def\csname PaperRevision@01s_graphic_equalisation_statement\endcsname{294}
\expandafter\def\csname PaperRevision@02_seldon_manifolds_theory_referenced\endcsname{31}
\expandafter\def\csname PaperRevision@02a_seldon_manifolds_psychological_grounding\endcsname{15}
\expandafter\def\csname PaperRevision@02b_seldon_manifolds_social_and_psychological_application\endcsname{7}
\expandafter\def\csname PaperRevision@02c_seldon_manifolds_decision_making\endcsname{11}
\expandafter\def\csname PaperRevision@03_law_of_narrative_momentum_theory_referenced\endcsname{32}
\expandafter\def\csname PaperRevision@04_narrativium_theory_referenced\endcsname{23}
\expandafter\def\csname PaperRevision@05_frames_relation_and_the_empty_limit\endcsname{11}
\expandafter\def\csname PaperRevision@06_assessing_an_account_of_this_kind\endcsname{29}
\expandafter\def\csname PaperRevision@07_causal_attitude_networks_on_the_substrate\endcsname{4}
\expandafter\def\csname PaperRevision@08_certified_answers_as_an_audit\endcsname{4}
\expandafter\def\csname PaperRevision@09_adaptability\endcsname{94}
\expandafter\def\csname PaperRevision@09a_adaptability_implementation\endcsname{7}
\expandafter\def\csname PaperRevision@10_applied_graphic_equalisation\endcsname{7}
\expandafter\def\csname PaperRevision@10a_red_queen_race_civilisational_rate_of_change_llms_biological_timescales_anthropocene_transition\endcsname{13}
\expandafter\def\csname PaperRevision@11_relational_information_ethics\endcsname{16}
\expandafter\def\csname PaperRevision@A1_constructive_graphic_equalisation_implementation_referenced\endcsname{8}
\expandafter\def\csname PaperGenerated@01s_graphic_equalisation_statement\endcsname{1}
\edef\PaperMetadataStem{\jobname}
\expandafter\ifx\csname PaperRevision@\PaperMetadataStem\endcsname\relax
  \PackageError{paper_metadata}{No metadata for \PaperMetadataStem}{Add the paper stem to paper_metadata.json.}
  \def\PaperRevision{?}
\else
  \expandafter\let\expandafter\PaperRevision\csname PaperRevision@\PaperMetadataStem\endcsname
\fi
\newcommand{\PaperDateLine}{\textbf{Draft \texorpdfstring{\textemdash}{---}\ \PaperDraft}\\[2pt]\PaperDisplayDate\ \texorpdfstring{\textbullet}{\textperiodcentered}\ document revision \PaperRevision}
\newcommand{\PaperStatementDateLine}{\textbf{Draft \texorpdfstring{\textemdash}{---}\ \PaperDraft}\\[2pt]\PaperDisplayDate\ \texorpdfstring{\textbullet}{\textperiodcentered}\ statement}
\expandafter\ifx\csname PaperGenerated@\PaperMetadataStem\endcsname\relax
  \newcommand{\PaperFooterStamp}{\textsc{\PaperDraft\,r\PaperRevision}}
\else
  \newcommand{\PaperFooterStamp}{\textsc{\PaperDraft\ / statement / from long r\PaperRevision}}
\fi
\newcommand{\PaperNoticeStamp}{\textbf{DRAFT --- \PaperDraft\ r\PaperRevision} (\PaperDisplayDate).}



\title{The Law of Graphic Equalisation\\\large Statement}
\author{Matthew Parslow\\Independent Researcher}
\date{\PaperStatementDateLine}

\begin{document}
\maketitle
\begin{center}
\fbox{\parbox{0.92\textwidth}{\small
\textbf{STATEMENT.} This is the theory; the companion \emph{Law of Graphic
Equalisation} (Theory) is the case for it. What is missing here is not detail but
\emph{argument} --- the removals that establish how little is required are the long
paper's bulk and its actual work, so anything below that a reader wants shown rather
than asserted is shown there. Calling this a short form would describe
it by its length instead of by what it does. Every labelled block below is pulled from the long paper at build time, and so are the
falsifiability commitments attached to them, so none of those can drift from it; the
commitments the long paper carries for material \emph{not} stated here are listed by
title at the end, so the omission stays checkable in this document rather than only in
the build log. The connecting prose is written here and is not synchronised, so it can.
Four companions are named in the material below without being carried by it: \emph{Is}
(the base claim this paper begins after), the physics companion, the mind companion,
and \emph{Narrativium}.}}
\end{center}
\vspace{1em}

\begin{abstract}
One minimal structure is proposed as underlying phenomena usually studied apart. This statement gives
the structure, the frame machinery through which anything reaches it, what happens when a write is
refused, the single prohibition it carries, and the commitments on which it fails. It argues for none
of them. The removals that establish how little is required are the long paper's bulk, and a reader
who doubts a particular claim here should expect to be persuaded there and not by this.
\end{abstract}

\section{The substrate}
The substrate is structure, the transformations operative upon it, and a grammar of which transformations
are admissible. Only the first is required for there to be anything; the other two are what a world in
which anything \emph{happens} costs. Nothing beyond the three is claimed primitive, and anything else exists
where a graph happens to represent it.

\begin{definition}[Graph]\label{def:graph}
A graph is
\[\boxed{G=(g,r,\Gamma),}\]
where $g$ is represented state, $r$ is operative transformation structure, and $\Gamma$ is the admissibility grammar governing graph-originated transformations.

Nothing in that is typed, and the omission is the definition's content rather than its brevity. What
the word \emph{graph} denotes here is the floor --- difference, and what differing consists in --- while
the typed object the word ordinarily calls to mind is derived and arrives with anything that runs
(Remark~3.11 of the long paper (\emph{What plurality forces, and what it does not})). A reader supplies the typing regardless, which is why this is said here
and again where it matters rather than left to the definition to carry alone.
\end{definition}
\begin{definition}[Grounding]\label{def:grounding}
A \emph{grounding} is a particular assignment of rule content and initial conditions to the substrate
--- one $(g,r,\Gamma)$ with its laws and its starting configuration fixed --- together with the
domain it is taken to model. The substrate is what every grounding shares; a grounding is what any
one of them adds. Claims made about the substrate are accordingly claims about all groundings, and
are correspondingly weak; claims made within a grounding carry that grounding's contingent rule
content and are where this account's falsifiable commitments mostly sit.
\end{definition}
\begin{definition}[Admissible continuation]\label{def:admissible}
A continuation $\delta$ belongs to $\Poss(G_t)$ iff it is expressible in the current graph vocabulary, satisfies the relevant antecedent conditions, and is admissible under the operative grammar for the transformation class in question.
\end{definition}
\begin{definition}[Collapse]\label{def:collapse}
Collapse is the frame-relative passage from a pre-realisation possibility structure to realised continuation:
\[\boxed{\Coll_t:\Poss(G_t)\rightsquigarrow G_{t+1}.}\]
It is not identified with deliberative choice.
\end{definition}

The Law is that a graph's admissible continuations are fixed by its present structure, and that
realisation selects one of them, which then becomes the antecedent for what follows.

\begin{schema}[Law of Graphic Equalisation]\label{hyp:the-law}
For realised graph $G_t$, its present structure specifies admissible
continuations $\Poss(G_t)$. Interaction and collapse realise $G_{t+1}$,
which becomes the antecedent for the next local transition. Here
\emph{equalisation} names this accommodation of a proposed continuation to
represented admissibility and containing constraints. It imposes no further
numerical balance or fixed-point condition beyond the stated recurrence.
\end{schema}

A realised deformation is not submitted for entry, so what a containing grammar declines is not undone
by the declining; it stays where it is. That remainder is the residual, and admitting it is how a grammar grows from residual.

\begin{definition}[Residual]\label{def:residual}
The residual is
\[\boxed{\epsilon=\delta^R\ominus\delta^P,}\]
where $\ominus$ is the grounding's difference relation preserving relevant discrepancy.
\end{definition}
\begin{proposition}[Residual is what a containing grammar does not admit]\label{prop:residual-inadmissible}
Because the graph is a subgraph of its containing frame $W$, a realised graph-originated deformation is \emph{already} a change in $W$ --- it is not submitted to $W$ for entry, and there is no step at which $W$ grants or withholds it. The delta is not offered; it \emph{is}. What the \emph{change grammar} of $W$ --- write it $\Gamma_{W_c}$, the change face of $W$'s grammar (Proposition~\ref{prop:two-gates}) --- admits coheres with the rest of $W$ and merges into its structure; what $\Gamma_{W_c}$ does not admit is not thereby undone, but remains in place, present and unintegrated. Admissibility sorts a change that already exists; it does not gate one seeking to occur. The residual is exactly this non-admitted remainder:
\[
\boxed{\epsilon \;=\; \delta^R \ominus \delta^P \;=\; \text{the part of the realised deformation }\Gamma_{W_c}\text{ does not admit.}}
\]
\end{proposition}
\begin{corollary}[Growth is grammar extension]\label{cor:mint-is-grammar}
A residual is absorbed either by a containing frame extending $\Gamma_{W_c}$ to admit it, or by paying to copy it out (Section~7 of the long paper (\emph{The frame interface and its economics})); otherwise it persists in place and is re-encountered each time the merge runs. The first route is the substrate form of \emph{minting} a new distinction: the residual does not fit because it is inadmissible, so admitting it requires extending the grammar, and extending the grammar is the act by which a previously inadmissible deformation becomes recognisable. Recognition and growth are therefore one mechanism joined at one place --- the residual is where the shape a graph decodes and the grammar that admits it meet.
\end{corollary}
\begin{remark}[Only $g$ is required; the triple is what a \emph{running} world costs]\label{rem:only-g}
The triple is not the minimum. Strip $r$ and
$\Gamma$ and something remains: structure, with relations among its parts. Nothing \emph{happens} in
it --- no transformation, so no collapse, no residual, and nothing it individuates --- but it has not
thereby become nothing, and it has not stopped being individuable by something else. What $r$ and
$\Gamma$ buy is \emph{occurrence}, and everything this account is about --- change, collapse,
observation, cost, experience --- requires them. $\Gamma$ is definitional --- an everything-admitting
grammar is still the grammar $\Poss$ is written over (Definition~\ref{def:admissible}) --- and what
Remark~\ref{rem:nothing-compulsory} calls an optimisation is the \emph{restrictiveness} of its two
faces, not its existence. The triple is the price of a world in which anything
happens rather than the price of there being anything, and Remark~\ref{rem:nothing-compulsory}'s
``only that rules run is compulsory'' should be read in that narrower sense.

A distinction is doing the work there. Being a
frame is being individuated (Definition~\ref{def:frame}), and individuation is performed by the
container, not by the thing: a boundary belongs to the frame that draws it. So a structure with no
rules of its own is still a frame --- something else instantiated its boundary --- and what it lacks
is not framehood but the capacity to individuate anything itself. Stripping $r$ does not leave ``no frame'', and reading it so would invalidate the mind companion's result, which reasons from a rock's being a frame to the universe's being one (its Proposition~\ref{MS-prop:minimal-universe}).
\end{remark}

What this paper postulates over \emph{Is} is the ledger of Table~\ref{tab:ledger}: named entries, not
a count, each with where it is argued. The prose defending each entry is the long paper's.

\begin{table}[htbp]\centering\small
\caption{The ledger: what this paper postulates over \emph{Is}.}\label{tab:ledger}
\begin{tabular}{@{}p{0.17\textwidth}p{0.36\textwidth}p{0.38\textwidth}@{}}
\toprule
\textbf{Status} & \textbf{Entry} & \textbf{Where argued, and what it replaced}\\
\midrule
Inherited &
Both of \emph{Is}'s propositions: to be is to stand in relation, and to cohere in it &
\emph{Is}, Propositions~\ref{IS-prop:is} and~\ref{IS-prop:cohere}. The second is its stated structural coherence condition; the informational rendering agrees with it but does not ground it (Remark~3.10 of the long paper (\emph{Coherence, not relation, is the criterion}))\\[2pt]
Admitted &
That anything \emph{happens} &
Unpaid, and one of two --- the other being \emph{Is}'s claim that there is anything at all\\[2pt]
Taken to follow &
The transformations $r$, and $\Gamma$ with them &
Something has to do the happening; the least graph that does anything is a rule. Two candidates remain open\\[2pt]
Postulated &
A realisation operator; a destroy\slash preserve distinction &
The attempt to dissolve collapse traded it for union\\[2pt]
Imported &
An information measure on represented relations &
Shannon-style measurement requires declared variables and distributions; absent at the base of \emph{Is}\\[2pt]
Operational consequence and methodological guard &
Axiom~\ref{ax:no-view} &
The weak no-informative-view result uses the measure (Remark~2.5 of the long paper (\emph{The axiom in bits})); the full frame-indexing discipline is stipulated here\\
\bottomrule
\end{tabular}
\end{table}

\section{Determinism and sufficiency}
What the Law requires of a transition and what it does not require of the substrate are separate
claims, and both are exposed here: an antecedent has to be there already for a continuation to be
selected from it, the selection adds no primitive randomness of its own, and nothing beyond the three
is needed to carry any of it.

\begin{proposition}[Antecedent-state principle]\label{prop:antecedent-state}
The conditions under which a transition is admissible are functions of its antecedent state. A consequence that exists only because a transition has occurred cannot retroactively serve as an antecedent cause of that same transition.
\end{proposition}
\begin{proposition}[Graph-law determinism]\label{prop:determinism}
For a fully specified antecedent configuration---including the graph state, operative transformation structure, admissibility grammar, and every causally operative containing relation relevant to the collapse---the realised continuation is determined:
\[
\boxed{
\left(G_t,\mathcal R_t^{\mathrm{causal}}\right)
\mapsto
G_{t+1}.
}
\]
The Law (Schema~\ref{hyp:the-law}) contains no primitive random actualisation operator.
\end{proposition}
\begin{hypothesis}[Minimal graph sufficiency]\label{hyp:minimal-graph}
No parent pointer, world model, affordance catalogue, Narrative, Prospect, Value, or self-model is
primitive; such structures may exist in $g$ when the graph actually represents them. Any additional
structure required only by a particular capability class is represented content or grounding, not a
primitive of graphhood, unless a graph lacking it cannot participate in the basic realised transition
relation.
\end{hypothesis}
\begin{hypothesis}[Discrete sufficiency]\label{hyp:discrete}
Every quantity required internally by the present substrate can be represented by discrete distinguishable states at the operative frame and resolution.
\end{hypothesis}

\section{Frames}
A frame is how anything is reached at all: there is no view from nowhere, so every claim is a claim
made in some frame, and the frame is not a container holding contents but the relation that
demarcates them.

\setcounter{axiom}{-2}
\begin{axiom}[Relativity of grounding]\label{ax:no-view}
There is no operational view from nowhere. Every distinction, measurement, identity relation, admissibility judgement, and represented quantity is made in some frame (Definition~\ref{def:frame}). Frames may represent one another, but no quantity used by the theory requires a frame-independent evaluator.
\end{axiom}
\setcounter{axiom}{0}

\begin{definition}[Frame]\label{def:frame}
A \emph{frame} is an instantiated boundary relation --- the edge $\mathcal B_{W:G}$ itself, and not the subgraph $G$ it demarcates. The subgraph is the relatum; the frame is the relation. To individuate is to instantiate such an edge, which is why individuation is per frame rather than a fact about the graph.
\end{definition}
\begin{corollary}[Frame and dimension are one edge type in two roles]\label{cor:frame-dimension}
A dimension is the same construction read from the other end. The boundary $\mathcal B_{W:G}$ is a dimension of $W$ --- one of the properties along which $W$ separates what it contains --- and each dimension $W$ instantiates is a frame to the values it separates. So \emph{frame} and \emph{dimension} name two roles one edge type plays relative to another edge, not two kinds of thing, and the asymmetry between them is positional in the same way that the asymmetry between a parent and a child is.
\end{corollary}
\begin{proposition}[One interface, two gates]\label{prop:two-gates}
A situated graph meets its containing frame through a single relation, and what an observation from that frame delivers is not distinct in kind from what the containing frame imposes: to a subgraph, being observed \emph{is} being deformed by its container. That one relation is gated twice, by the same kind of admissibility grammar on its two faces --- the \emph{aperture grammar} $\Gamma_a$ on the input and the \emph{change grammar} $\Gamma_c$ on the output:
\[
\boxed{\underbrace{\text{aperture}}_{\Gamma_a}\;\longrightarrow\;\underbrace{\text{rules}}_{r}\;\longrightarrow\;\underbrace{\text{change}}_{\Gamma_c}.}
\]
$\Gamma_a$ limits, by grain, which of the container's deformation enters the rules; $\Gamma_c$ determines which of the rules' realised deltas cohere into the container and which remain as residual (Proposition~\ref{prop:residual-inadmissible}). Both are faces of the frame's one grammar $\Gamma$, and instantiate for a specific container $W$ as $\Gamma_{W_a}$ and $\Gamma_{W_c}$; it is $\Gamma_{W_c}$, the change face, that the residual proposition uses. Execution runs between the two.
\end{proposition}
\begin{definition}[Aperture]\label{def:aperture}
A frame's \emph{aperture} is its input gate $\Gamma_a$ together with what that gate fixes: which
dimensions it can separate along at all, how finely it separates within each (its resolution,
per dimension, for which this paper also uses \emph{grain}), and how far along each it reaches. The three are properties of one edge --- existence,
grain, and extent --- and are named separately because they vary independently, not because they are
different kinds of limit. What falls outside an aperture on any of the three is \emph{absent} to that
frame rather than held-but-unresolved.
\end{definition}
\begin{definition}[Change grammar]\label{def:change-grammar}
A frame's \emph{change grammar} $\Gamma_c$ is its output gate: it determines which of the realised
deltas its rules produce cohere into the containing frame and which do not. What it declines has
entered the frame and been operated on, so it is \emph{residual} (Definition~\ref{def:residual})
rather than absent, and it persists in place rather than being discarded. $\Gamma_a$ and $\Gamma_c$
are the two faces of one grammar $\Gamma$ on one boundary edge, and the asymmetry between them is
which direction they gate, not what kind of thing they are.
\end{definition}
\begin{definition}[States of a dimension]\label{def:states}
For a frame $\mathcal F$, a dimension is in exactly one of four states:
\[
\boxed{\mathrm{NULL}\;\neq\;\mathrm{undetermined}\;\neq\;\mathrm{undisclosed}\;\neq\;\mathrm{resolved}.}
\]
\emph{NULL}: the dimension is not instantiated; there is nothing to have a value. \emph{Undetermined}: the dimension is instantiated and \emph{no value has been realised} --- collapse has not occurred. \emph{Undisclosed}: a value has been realised but is not held by this frame. \emph{Resolved}: a realised value, held.
\end{definition}

Provenance is a relation to what produced a structure or event. Its attribution
is inferred from evidence available in a frame; an origin record is evidence,
not a self-authenticating intrinsic mark. A local attribution can be unique
under a declared model. This is distinct from assumption-free closure of the
recursively supplying operators, which is separately argued in the long paper.

\begin{definition}[Relational provenance]\label{def:provenance}
Provenance is a relation between a structure or event and the causal
organisation or history that produced it. A provenance attribution is a
frame's account of that relation, constrained by the evidence the frame
can distinguish. Neither an origin nor a source is an additional intrinsic
mark carried independently of the relations and traces that make it
distinguishable.
\end{definition}
\begin{remark}[Status of the principle]\label{rem:provenance-status}
This is a derivational treatment of an established concept. Source-monitoring
theory treats remembered origins as judgments based on memory characteristics
and other available information (Johnson et~al., 1993). Causal identification
characterises when causal queries can be recovered from observational or
interventional information under stated models (Shpitser and Pearl, 2008).
Database provenance formalises relationships between derived outputs,
contributing inputs, and source locations (Buneman et~al., 2001). These are
distinct precedents with their own scopes; none is being claimed as a
derivation of this entire ontology.

\emph{Provenance is inferred, not inherent} states the relational discipline
for origin. Actual causal relations do not wait for an attribution. What a
frame knows of them comes through its relations to present structure,
records, and effects, under Axiom~\ref{ax:no-view}. The principle adds no
origin primitive or further observer to the substrate. The attribution is
itself represented structure and can affect subsequent action.

A provenance record can be carried explicitly. Such a record is evidence
whose relation to its claimed source must be established in the frame;
writing an origin into a label does not make the label authenticate its own
history. Conversely, reliable recording and controlled intervention can
make a local attribution determinate. The principle does not say that
every origin is unknowable or that actual histories are interchangeable.

Graphic Equalisation's contribution here is derivational backing from its
relational commitments and a common account of the principle's applications.
The dependencies are explicit: the relational ontology supplies the account
of origin as a causal relation; relative grounding places its operational
characterisation in a frame; the declared evidence map then yields the
recovery or non-recovery result. The no-view axiom alone does not prove the
absence of an intrinsic historical property. The backing is conditional on
these foundational commitments and the distinction between actual relation
and represented attribution.
The evidence-relative proposition below states a standard identifiability
argument in that account's terms; neither the underlying concept nor the
argument is claimed as a new discovery. ``Inferred'' also does not require
conscious deliberation: source-monitoring judgments can be automatic or
deliberate (Johnson et~al., 1993).
\end{remark}
\begin{proposition}[Provenance recovery is evidence-relative]\label{prop:provenance-inference}
For a declared family of candidate histories \(H_F\), let \(O_F(h)\) be the
evidence distinguishable in frame \(F\), and let \(\pi_F(h)\) give the
provenance distinction being asked about. The compatible provenance set is
\[
P_F(s)=\{\pi_F(h):h\in H_F,\ O_F(h)=s\}.
\]
If two compatible histories have different provenance, the evidence \(s\)
does not determine a unique attribution. A singleton \(P_F(s)\) determines
that distinction only relative to the declared history family and evidence
model. An empty set indicates that the model supplies no compatible history.
\end{proposition}

What a frame's gates admit, and what it can reach from where it stands, are two quantities and not
one; a disclosure that moves a gate changes what the gates admit, and that is where an information hazard is
located here rather than treated as a class of its own.

\begin{definition}[Capability and accessibility]\label{def:capability}
A frame's \emph{capability} is what its gates admit --- what it can represent and what it can do. Its
\emph{accessibility} is what it can \emph{reach} from where it stands, given what it holds to spend. A
disclosure moves a gate and changes capability; whether the reachable point can then be crossed to
depends on the Narrativium measured for the frame. \emph{Narrativium}, which defines that quantity, types it as a
\emph{capacity} rather than as a cost (its Remark~\ref{NV-rem:narrativium-standing}), and keeps measurement and expenditure distinct; the phrasing
here follows it.
\end{definition}
\begin{definition}[Information hazard]\label{def:info-hazard}
A disclosure changes admissibility when it moves a frame's aperture gate, its change grammar gate, or
both (Proposition~\ref{prop:two-gates}), and what it changes when it moves one is that frame's
\emph{capability} (Definition~\ref{def:capability}). An \emph{information hazard} is a capability
change whose negative sign is admissible in the reference frame, above the point at which the impact
forces a write. The class is disclosures that change admissibility; the sign is applied afterwards by
a frame with preferences, and the magnitude is a ratio to what the frame can absorb rather than an
absolute.
\end{definition}
\begin{proposition}[The threshold is a ratio, and sits where the impact forces a write]\label{prop:hazard-threshold}
A minimum magnitude is required, or every capability change would be a hazard. The threshold cannot be
an absolute, since an absolute would be a magnitude fixed independently of any frame, which
Axiom~\ref{ax:no-view} forbids; it is a ratio --- impact against the scale of the frame bearing it.
And the substrate supplies where the ratio sits rather than leaving it to stipulation:
\emph{the threshold is the point at which the impact forces a write}. An impact a frame can absorb
coheres with what it holds, forces no write, and costs it nothing; an impact it cannot absorb leaves
residual and compels change. That is the same line every other quantity here is measured against, and
it is a ratio by construction, because absorption capacity is the frame's own.
\end{proposition}

\section{Identity}
Sameness is not read off the substrate. It is fixed by whatever frame individuates, which leaves two
identity relations where ordinary usage has one --- the graph's own continuity with what produced it,
and the tracking a container does of a situated instance.

\begin{corollary}[No frame-free SUPERGRAPH tracking identity]\label{thm:no-frame-free}
The substrate alone supplies no $\operatorname{Id}_W$ tracking edge without
an individuating frame $W$. This follows from Definitions~\ref{def:frame}
and~\ref{def:supergraph-identity}; it does not deny SELF continuity across an
actual run (Definition~\ref{def:self-identity}).
\end{corollary}
\begin{definition}[SELF identity]\label{def:self-identity}
SELF identity is the continuity of the graph as a whole --- its rules together with the state they
act on --- across realised transitions. A graph is the same self while its present remains
congruous with its own history: if its rules changed abruptly, the structure carried forward would
no longer cohere with what produced it, and that incongruity is the discontinuity. Represented
self-description --- content inside $g$ by which a graph identifies itself --- is one optional part
of this and not the whole of it; a graph may have SELF identity without representing itself at all.
\end{definition}
\begin{definition}[SUPERGRAPH identity]\label{def:supergraph-identity}
For graph $G$ individuated in frame $W$,
\[\iota_{W:G}=\operatorname{Id}_W(G)\]
is the identity relation under which $W$ tracks that situated instance.
\end{definition}

\section{Economics}
Cost characteristically attaches to writes rather than to holding, and what a write costs is set by
whatever runs the rules rather than by any universal. A graph changes itself directly and its
container only as the region of it that it is, so the distribution of what is free, what is costly and
what is directly writeable is a characterisation of a common configuration rather than a necessity.

\begin{proposition}[Cost characteristically reduces to the write]\label{prop:economics}
In the substrate as characterised here, the boundary relation itself is free, and the standing structure persisting is free \emph{to the frame}; the operation that costs is a delta that writes over that structure, and $r$ is what writes. Whether persistence is free \emph{simpliciter} is set by the rules of the superframes, which may sustain a frame without charge or require a write to keep it standing, and the substrate does not decide between them. \emph{Free} to the frame means unbilled to that frame. Information preservation is a separate property: evaluating a rule can be a reversible write, but reversibility alone does not fix its physical cost (Remark~7.17 of the long paper (\emph{Where the check goes, and why the free case is the empty one})). To that extent
\[
\boxed{\text{cost is paid by rules.}}
\]
\end{proposition}
\begin{proposition}[Self-change is direct and gated; a graph is a subset of its container]\label{prop:no-direct-self-alter}
A graph changes itself and its subgraphs directly. Self-change is real, but it is rule-governed, with its own cost and limits (an instance of the non-uniform writeability of Remark~\ref{rem:writeability}); a substrate's plasticity is exactly the rules for how its rules may change. What a graph does \emph{not} do is change its container directly: the container applies deformations to it --- being observed \emph{is} being so deformed --- and gates its changes through its own change grammar (Proposition~\ref{prop:residual-inadmissible}). But a situated graph is a \emph{subset} of its container, not a thing set against it, so a self-change is already a change in the container, in the region that is the graph, propagating outward through whatever couples it to the rest.
\end{proposition}
\begin{remark}[Writeability status is contingent, not uniform]\label{rem:writeability}
The write-economics stated here --- reads free, cost carried by writes, and no direct alteration of the container (Proposition~\ref{prop:no-direct-self-alter}) --- should be held at the strength the evidence supports, and the two halves of that need
separating. The \emph{structural} claim --- that a graph affects its container as a subset of it and not
by some other route --- is a strict restriction and is committed to as one
(Section~19 of the long paper (\emph{Falsifiability commitments})). What is a pattern rather than a universal is the \emph{distribution}
below: which writes are free, which are costly, and how much is direct. What is freely writeable, what carries a cost, and what a graph can write directly rather than only across a boundary, need not be uniform; different structures may have different writeability profiles. The strict forms above characterise a common and efficient configuration (Remark~\ref{rem:nothing-compulsory}) rather than a necessity of graphhood, and asserting a single strict write-economics everywhere would be the same over-specification Remark~\ref{rem:nothing-compulsory} warns against, a fixed value put where the account has only a setting. These are accordingly carried as falsifiability commitments, open to counterexample, and a structure with a different writeability profile refines the characterisation rather than breaking the substrate.
\end{remark}
\begin{definition}[Well-formedness]\label{def:well-formed}
A graph is \emph{well-formed} to the degree that it pays for nothing that does no work for it. Four costs bear on that and they behave differently, so the notion is graded along each rather than binary.

\emph{Writes} are paid when made (Proposition~\ref{prop:economics}), which prices doing. \emph{Retention} is nearly free per unit, which is why an unused distinction is not immediately pruned. \emph{Search} is paid on every operation over what is held and scales with how much is held, so inert structure is cheap to keep and expensive to keep in quantity.

The fourth is the worst and is not a variant of the third. A write may land on structure that cannot affect anything --- a modification to a branch nothing reaches. It is paid in full: decided, made, checked, risked. It returns nothing. And unlike a search, which is expensive and then yields, this failure is \emph{silent}: the graph has spent, has no effect to show for it, and no signal that the effect is missing. What it acquires instead is a false model of itself, in which something was changed that was not --- residual it cannot see, produced by its own inert contents, and the share of writes exposed to it grows with how much is held.
\end{definition}
\begin{proposition}[Evaluation is a write, and free because information-preserving]\label{prop:evaluation-writes}
A read is an operation that returns a datum; the returned datum has to be somewhere to be used, and in
a substrate that is graph all the way down there is no return channel outside the structure, so the
value materialises as structure and the write is the return. Evaluation is a write in that plain
sense: a match outcome, a position in an ordering, and the fact of having tried are returned values,
and returned values are structure. It is information-preserving where its return map is injective. The reference economics leaves it unbilled to the evaluating frame; neither statement establishes zero physical cost. What it costs does not follow from that: a rule's
cost is set by whatever supplies its primitives and is inherited along provision, so reversibility is a
structural fact about a map while cost is a fact about what supplies a primitive, and the two must not
be run together.
\end{proposition}
\begin{proposition}[Primitive leakage and the scope of provenance recovery]\label{prop:leak}
Primitive effects can disclose supplying structure; leakage is ordinary in
the cases of Remark~14.2 of the long paper (\emph{Exposure is excluded; leakage is produced}). A declared local provider map
can be identified where its compatible provenance set is a singleton.
Where two admissible supply histories produce identical available evidence
but differ in requested provider provenance, that distinction is not
identified (Proposition~\ref{prop:provenance-inference}).

Local identification under an evidence and supply model does not certify
that model's exhaustive coverage of an unrestricted supply chain. The
account's nonclosure of that unrestricted question is a separately argued
commitment of Section~15 of the long paper (\emph{The single prohibition, and its formal core}), not an unconditional consequence
of the evidence-relative proposition or of the fact that inference itself
uses primitives.
\end{proposition}

\section{When a write is refused}
An inadmissible write yields no continuation in that frame. What happens after that --- whether
anything propagates, whether a value is produced instead --- is rule content and not substrate. The
failure itself does not vary; the frames do, and a handler is an installed policy rather than a
discovered semantics.

\begin{proposition}[Decoherence is one operation; the frames are what differ]\label{prop:decoherence-one}
Decoherence is one operation. It does not come in kinds and it
does not come in degrees, and intended and unintended incoherence both return no coherent state. What differs between them is the
frame --- arithmetic in one, kinematics in the other --- and \emph{the frames are what make the cases
different, not the failure}. ``Everything in the superframe'' is a constant expression whose value is
frame-relative, which is Axiom~\ref{ax:no-view} applied to incoherence rather than to observation.
\end{proposition}
\begin{proposition}[A handler relocates a failure; it does not close it]\label{prop:handler-relocates}
A handler is installed policy that determines what a frame coheres a refused write to, and it can be
set to cohere it to anything, including success. Swallowing does not undo the inadmissibility: the
containing grammar still refused and the write still did not happen, so what is removed is the frame's
\emph{record} of the refusal rather than the refusal. Detection is available and never free --- a frame
holding a record of the attempt separately from the handler's report can compare the two --- but it is
unavailable from the handler's own report, which is self-consistent by construction. The audit is
therefore installed structure, and being itself a policy it is abusable in the same way one level in,
so what its purchase buys is relocation rather than closure, terminating only in a containing frame.
\end{proposition}
\begin{proposition}[Dependents are held by the referring frames, not the referent]\label{prop:vacated-three-states}
A thing does not carry a record of what depends on it: the dependencies live in the frames holding the
references and not in the referent, so from the position of the structure being removed they are not
visible at all, and knowing what would break requires the containing frames. Removal therefore yields
three states rather than two. \emph{Absent} --- nothing ever cohered, nothing is held and nothing owed,
and it is free. \emph{Vacated with the references left standing} --- a hole, unpaid, with consequences
determined by whatever the dangling relations do next. \emph{Vacated with the absence represented} ---
a tombstone, a cascade rule, a retained count of dependents --- which resolves cleanly and is installed
machinery paid for under Proposition~\ref{prop:economics}.
\end{proposition}
\begin{proposition}[Termination propagates along supply, not containment]\label{prop:termination-supply}
Being contained by something and being run by it are two relations, and they need not coincide.
Termination propagates along the supply relation and not the containment relation: downward it need not
reach at all, a container being able to stop while what it held runs on, and upward the supplier's stop
is total, at whatever containment distance and with no handler available, since the operations a
handler would run are the ones withdrawn. This splits the dependents of a stopped structure into three
kinds rather than two. A dependent that merely \emph{refers} to it is left dangling and keeps running,
wrongly; a dependent \emph{wholly run by} it simply stops; and a dependent only \emph{partly} supplied
by it keeps running with an incomplete operation set, and its behaviour is undefined.
\end{proposition}
\begin{proposition}[A boundary is known only by crossing it]\label{prop:boundary-by-crossing}
A frame that never fails discloses nothing about what it admits. Admissibility is invisible from inside
its own range, since everything works, and the boundary is legible only where it is crossed --- so the
way to learn a rule is to force the failure and read the residual. The test and the change are
accordingly the same operation: knowledge of a boundary is bought by crossing it, and there is no
cheaper instrument. It follows that every handler installed reduces what failure can teach, a frame
that coheres its failures to tidy values reporting success and disclosing nothing, so that robustness
and legibility are traded against one another rather than being two names for quality.
\end{proposition}

\section{Time}
Time is not a substrate dimension. It is a dimension a frame constructs from represented succession.
A closure cycle is the simplest self-contained, repeatable operational unit for that construction, and
is therefore the default hypothesis where no finer or externally supplied counter is represented. It
is not a substrate requirement: a frame cannot establish from inside whether an apparent cycle is
self-closing, clocked, superframe-scheduled, or realised at a fixed hidden duration.

\begin{axiom}[No substrate time dimension]\label{ax:no-time}
\emph{Time is not a first-class dimension of the substrate.} What the substrate has is \emph{change}: realised collapses. It has no ratchet over and above them --- no mechanism that manufactures succession --- because the three properties usually bundled under that word are readings over changes or properties of them, not a further thing the substrate runs. The ordinal event order of Section~11 of the long paper (\emph{Causal order and local time}) is carried by the changes' own enabling-relations; direction is the asymmetry of enablement itself --- if one change makes another possible, the second
does not thereby make the first possible, which holds whether or not a change preserves what it
followed; and persistence is only that further changes occur. A time dimension exists only where a frame instantiates one \emph{to represent} that order.
\end{axiom}
\begin{proposition}[Time is inferred, not observed]\label{prop:time-inferred}
No collapse presents a temporal coordinate. Succession is a relation \emph{between} realised states, and a relation between states is not a datum of any one of them. A frame therefore acquires a time dimension only by \emph{inferring} it from the relations among what it holds. Consequently
\[
\boxed{\text{a frame observes time}\;\Longrightarrow\;\text{that frame is capable of inference.}}
\]
\end{proposition}
\begin{proposition}[Closure is a provisional operational clock hypothesis]\label{prop:cycle-unit}
A frame may use a recurrence that it treats as a closed rule-loop as its temporal unit. This is the
simplest self-contained explanation of a repeatable succession, so it is the default where no
competing counter is represented. The resulting cycle-count is operationally useful and has no fixed
hidden duration: one apparent cycle may consume variable change-load, be clocked from outside, or be
scheduled irregularly by a containing frame. None of those alternatives is decidable from inside when
all accessible clocks are rescheduled with it.

Closure is therefore neither necessary nor established at the substrate. A frame can infer a temporal
coordinate from any represented succession or counter that is adequate for its task; a closure cycle
is one economical construction of such a coordinate. A one-shot computation has no \emph{self-generated}
recurrence, but may still be embedded in, or supplied with, a represented ordering by another frame.
The present is a current realised boundary, not proof that the rule which led to it closes.
\end{proposition}
\begin{definition}[The present, structurally]\label{def:present}
Within a frame, the present is the collapse boundary at which every instantiated \emph{undetermined} dimension relevant to that collapse has been determined:
\[
\boxed{U_{\mathcal F}=0.}
\]
$U_{\mathcal F}$ counts undetermined dimensions only. NULL dimensions are excluded --- they are not variables. \textbf{Undisclosed dimensions are also excluded}: they are already determined, and a frame's failure to hold them is an epistemic deficit rather than an unresolved collapse. A frame with outstanding undisclosed structure is still at its own present.
\end{definition}
\begin{proposition}[No implicit historical state]\label{prop:no-implicit-historical-state}
A previous success, failure, attempt, or residual has no later causal force merely because it occurred. It affects later dynamics only insofar as its consequences are present in later graph or boundary state.
\end{proposition}
\begin{corollary}[History principle, strengthened]\label{cor:history}
\[
\boxed{\text{history matters only where it survives as present structure}.}
\]
The past is surviving consequence instantiated now, not a stack of still-existing historical frames.
\end{corollary}

\section{Primitives, and the single prohibition}
\begin{definition}[Primitives]\label{def:primitives}
A graph's \emph{primitives} are the operations it invokes without performing: the operator set its
rules bottom out in. Every level has them, and having them is what being a level consists of --- a
graph with no primitives would have to perform everything it invokes, which is the regress this
substrate does not contain. That is the whole of the definition, and in particular it says nothing
about what a graph can learn of where they come from.
\end{definition}
\begin{remark}[What is required and what is optimisation]\label{rem:nothing-compulsory}
Only that rules run is compulsory for anything to \emph{occur} --- structure alone is what is required for there to be anything at all (Remark~\ref{rem:only-g}). Applying a delta is itself a rule, not a primitive operation beneath the rules, so the specific dynamics --- inertia, the merge behaviour, closure --- are contingent rule content, and a different $r$ or $\Gamma$ is a different physics rather than a violated law (this is why the initial and boundary conditions of a grounding are supplied separately from its laws; Section~11 of the long paper (\emph{Causal order and local time})). The two gates of Proposition~\ref{prop:two-gates}, and even the retained standing structure itself, are \emph{optimisations} on bare rule-running: the aperture grammar $\Gamma_a$ bounds the input, the change grammar $\Gamma_c$ sorts the realised output into what coheres and what is residual, and the retained structure buys free continuation. A substrate whose grammar admits everything still runs --- effectively ungated, without momentum, wastefully. It is not a substrate \emph{lacking} $\Gamma$, since Definition~\ref{def:admissible} defines the possibility structure through the grammar and a collapse needs something to select from. The two gates are accordingly what an efficient graph accretes under a budget, not constituents of graphhood. The retained structure specifically is what distinguishes a self-sustaining loop from a purely reactive one, since coasting is the cached structure.
\end{remark}
\begin{proposition}[Primitives are supplied, not intrinsic]\label{prop:rule-runner}
A graph is data --- nodes, relations, and rules --- and the operations its rules bottom out in are supplied to it rather than contained in it. The same graph may therefore be run against different supplies of the same operator set, and nothing supplying an operator need be the substrate the graph describes or is ``in.'' Where the work is done is a fact about that supply and not about the graph.
\end{proposition}
\begin{proposition}[A well-formed runner does not expose itself, and no level is self-transparent]\label{prop:runner-opaque}
The asynchronous, locally synchronous structure a frame observes (Remark~11.3 of the long paper (\emph{Globally asynchronous, locally synchronous})) is the graph's \emph{internal} view and does not fix its runner's actual mode. Where a frame has no boundary to the substrate that runs it, its runner lies beyond the frame's epistemic boundary (Section~1 of the long paper (\emph{Scope and epistemic status})); a hidden runner may execute the graph serially while the graph appears asynchronous to itself (Remark~11.3 of the long paper (\emph{Globally asynchronous, locally synchronous})). And no level is self-transparent: reading is an operation whose product lands one level up (Remark~7.17 of the long paper (\emph{Where the check goes, and why the free case is the empty one})), so a level cannot hold the reading of itself while being the level that is running --- a frame may be given sight of what runs it and cannot be given sight of its own running.
\end{proposition}
\begin{proposition}[A primitive is another graph's running, exposed]\label{prop:runner-recursion}
Nothing supplying a primitive is a substance apart from graph data. An operator is the observable running of some other graph, which has primitives of its own, supplied in the same way. A graph's ``primitives'' are therefore the observable running of other graphs, exposed as an operator set, with no primitive runner at the base. The regress bottoms not at a floor but at the epistemic frontier of Section~1 of the long paper (\emph{Scope and epistemic status}): below the lowest measured level is the unclaimed pre-relational substrate.
\end{proposition}

Stepping back: the account forbids very little, and what it does forbid is one thing recurring rather than a list.

\begin{proposition}[The single prohibition]\label{prop:prohibition}
The account forbids one thing: \emph{a frame cannot get outside itself}. Concretely, and confined to
what is actually derivable here, it takes three forms. A frame cannot be transparent to itself at the
level it is running at (Proposition~\ref{prop:runner-opaque}). It cannot measure its own throughput in
units independent of it, since the measurement is funded by the throughput --- which leaves a number
obtainable and its absolute scale not. And it cannot hold a view from nowhere, which is Axiom~\ref{ax:no-view}.
\end{proposition}

That prohibition has a formal analogue --- under classical bivalence, a language strong enough to represent its own syntax cannot define its own
truth predicate, and must be given it from a metalanguage --- and the analogue is proved, though the
mapping from frames to formal systems is argued rather than proved. What is forbidden is complete
self-capture and not partial self-modelling, which is why reflective towers work and a frame can know
a great deal about itself. Two candidates are deliberately not on the list: that a frame cannot obtain
a quantity whose unit does not close within it is contingent rule content, and that a frame cannot
decide anything without writing is a pattern of the write-economics rather than a universal.

\section{Carvings}
Since the rules draw the boundary and the rules are contingent, no carving is privileged. That does
not make the notion vacuous: an ill-formed cut costs more writes, so a frame under a budget tends to
be running cheaper cuts, and what looks like a joint is where the cheap cuts of frames doing shared
work coincide. The step from shared work to similar cuts rests on a posit, labelled here as one.

\begin{hypothesis}[Near-uniqueness of the cheap cut]\label{hyp:cheap-cut}
For a given task under a given binding constraint, the cheapest arrangement that performs it is nearly
unique. This is an empirical posit and not a result of this account: nothing here derives it, and what
is available without it is weaker --- a binding constraint narrows the solution space, and the more
binding it is the narrower it gets.
\end{hypothesis}
\begin{proposition}[Convergence is observed, not enforced]\label{prop:convergence-observed}
Two frames doing the same work converge on similar cuts as each becomes better formed, because the
cheap cut for a given task is nearly unique. This is a tendency of cost and not the discovery of a
privileged boundary, and the two are distinguishable in principle: a privileged boundary would
\emph{enforce} convergence, and enforcement would extend to the hard cases; cost-shaped cuts converge
where the work is shared and come apart where it is not. The step from shared work to similar cuts
assumes Hypothesis~\ref{hyp:cheap-cut}.
\end{proposition}
\begin{proposition}[Convergence under compression]\label{prop:compression-convergence}
A structure operating under a constraint, and under pressure to spend nothing on what does no work,
approaches the arrangements that constraint forces --- asymptotically, and only over the region the
constraint touches. Two such structures therefore come to resemble each other without either having
access to the other, because both are downstream of the same forcing rather than converging on a
target.
\end{proposition}
\begin{hypothesis}[Located degradation]\label{hyp:located-degradation}
Where a classification's purpose can be specified independently and in advance of inspecting its
failures, the failures concentrate near the boundary of that purpose more than a uniform
distribution over the domain would produce. The claim is comparative and is not that failures occur
nowhere else.
\end{hypothesis}

\section{Falsifiability}
The substrate risks little by design and the groundings carry the risk. Below are the commitments
attached to the material this statement carries. They are a subset: the long paper carries the full
set, including those for results not stated here, and the omitted commitments are listed below by
title, so the omission stays visible rather than becoming a claim.

\begin{enumerate}[leftmargin=*]
\item \textsc{[shape]}~\textbf{Time is inferred, not primitive.} Fails if a frame's realised dynamics require a temporal coordinate that is not constructed from relations among what it holds --- a collapse presenting a duration as a datum of the state rather than as a reading over states (Axiom~\ref{ax:no-time}, Proposition~\ref{prop:time-inferred}). A frame that represents succession is not the counterexample; that is what the claim predicts. Marked \textsc{[shape]} because what would fail is the range of everything derived through the absence of a substrate time dimension.
\item \textsc{[detail]}~\textbf{Closure is the default operational clock hypothesis.} The claim is not that closure is required, but that a frame with no contrary evidence should prefer a treated-as-closed recurrence because it is the shortest self-contained clock model. It is revised, not falsified, where a represented counter, an external clock, or a superframe schedule improves operational use (Proposition~\ref{prop:cycle-unit}).
\item \textsc{[shape]}~\textbf{Two gates, and absence is not residual.} Fails if what a frame's aperture excludes is nonetheless held or disturbs it --- an excluded input leaving a remainder in place --- so that the input and output faces stop differing in what their exclusions cost (Proposition~\ref{prop:two-gates}, Definitions~\ref{def:aperture} and~\ref{def:change-grammar}).
\item \textsc{[shape]}~\textbf{Partial-order sufficiency.} Fails if realised dynamics require a total order over collapses in distinct frames --- a global simultaneity --- rather than only the boundary-induced partial order.
\item \textsc{[shape]}~\textbf{Composition.} The arrangement claim needs a comparative test; ``equal facility'' alone has no
measurement rule. One candidate protocol fixes a code, lossless domain-to-graph
maps, permitted auxiliary machinery, reconstruction criterion, and resource
budget before rendering matched target and adversarial control domains. Let
$d(D)=L_{\rm GE}(D)-L_{\rm native}(D)$ include the mapping, added rules and
residual description lengths. The comparative prediction is
$\overline d_{\rm control}-\overline d_{\rm target}>0$; a nonpositive
difference challenges that advantage under the declared protocol.
Independent renderers and a prespecified uncertainty rule are needed for a noisy
sample. This protocol has not been registered or run. It tests compression
advantage under its fixed code, not the ontology, and no operational composition
falsifier is claimed to have been passed. Independent recurrence in a domain's
own idiom remains suggestive evidence, also compatible with shared constraints.
\item \textsc{[shape]}~\textbf{Minimal graph sufficiency.} Fails if basic realised graph dynamics require a primitive beyond $(g,r,\Gamma)$ rather than represented or relational structure (Hypothesis~\ref{hyp:minimal-graph}).
\item \textsc{[shape]}~\textbf{Discrete sufficiency.} Fails if a claimed observable consequence cannot be preserved by any discrete representation, including dynamically varying quantisation, but is preserved by a genuinely non-discrete primitive (Hypothesis~\ref{hyp:discrete}).
\item \textsc{[shape]}~\textbf{Antecedent-state principle.} Fails if a realised transition can require as antecedent a fact that exists only as consequence of that same transition, absent prior fixed-point structure (Proposition~\ref{prop:antecedent-state}).
\item \textsc{[shape]}~\textbf{Graph-law determinism.} Fails if two genuinely identical, causally complete antecedent configurations produce different realised successors without any additional causally operative distinction (Proposition~\ref{prop:determinism}). Variation observed only from a frame that omits determining dimensions does not constitute such a falsification.
\item \textsc{[shape]}~\textbf{No implicit historical state.} Fails if an unrepresented past event alters later transition structure despite leaving no causal trace in graph or boundary state.
\item \textsc{[shape]}~\textbf{Frame-relative identity.} Fails if sameness and persistence are derivable from $(g,r,\Gamma)$ alone without individuation.
\item \textsc{[detail]}~\textbf{Located degradation.} Fails if, for classifications whose purpose is specifiable in advance of inspecting their failures, failure cases are distributed no differently near the purpose boundary than across the rest of the domain (Hypothesis~\ref{hyp:located-degradation}).
\item \textsc{[shape]}~\textbf{Convergence is not enforced.} Fails if convergence between independently developed classifications of one domain extends to the cases where the classifications' purposes differ, which is what a privileged boundary would produce and cost-shaped convergence would not (Proposition~\ref{prop:convergence-observed}). Marked \textsc{[shape]} because a privileged boundary is not a mechanism this account could replace, since the contingency of the rules that draw a boundary is what Section~16 of the long paper (\emph{Carvings, and why classification degrades}) rests on and what Axiom~\ref{ax:no-view} asserts one level up. Exhibiting one costs range rather than a mechanism.
\item \textsc{[shape]}~\textbf{Provenance is never closed by inference.}
The claim concerns assumption-free final closure of the recursively supplying
operators, not identification of a declared finite provider map. It is
challenged by an in-frame procedure establishing exhaustive unrestricted
supply coverage without assuming the completeness of its candidate model.
Where admissible supply alternatives have identical available evidence but
different provider provenance, they remain unidentifiable
(Proposition~\ref{prop:provenance-inference}). Local identification from leakage
or reliable records is permitted (Proposition~\ref{prop:leak}). The unrestricted
nonclosure is the separately argued frontier commitment of
Section~15 of the long paper (\emph{The single prohibition, and its formal core}), not a theorem derived merely from evidence dependence.
\item \textsc{[detail]}~\textbf{Solutions are nearly unique under a binding constraint.} Hypothesis~\ref{hyp:cheap-cut}, the posit under Propositions~\ref{prop:convergence-observed} and~\ref{prop:compression-convergence}. Fails if structures optimising hard against the same binding constraint, over the same region, are exhibited arriving at arrangements that do not resemble each other --- in which case constraint narrows the space without selecting within it, and both propositions weaken to that.
\item \textsc{[detail]}~\textbf{Cost is the write.} Fails if the boundary relation itself, or the retention of standing structure, carries an irreducible cost that is not any write's --- a frame charged for holding its boundary open, or for keeping structure standing, where no delta is written and none is occasioned. Those two are the non-write costs Proposition~\ref{prop:economics} itself names as free, so they are where the claim is exposed; reading is not the test, since a paid read is a write here (Proposition~\ref{prop:evaluation-writes}) and the falsifier would be absorbed rather than met. Marked \textsc{[detail]} because Remark~\ref{rem:writeability} already carries the distribution of costs as a pattern rather than a universal: a structure with a different writeability profile refines the characterisation rather than breaking the substrate.
\item \textsc{[shape]}~\textbf{Self-change is subset-change.} The claim is that a graph changes itself and its subgraphs directly (rule-governed), and affects its container only as a subset of it, by propagation through coupling. Fails if a graph changes its container by a means other than changing the region of it that the graph is --- a container-change not propagating from the graph's own subset --- or if a \emph{container-supplied} graph can self-terminate without acting on the container (Corollary~10.7 of the long paper (\emph{Severing the channel removes termination})); a self-supplied graph stopping itself is not a counterexample, being the case that corollary excludes.
\item \textsc{[shape]}~\textbf{Minting is a rule.} Fails if grammar-extension (growth) requires a primitive operation beneath the rules rather than a reflective rule subject to the same economics, admissibility, and antecedent-state principle (Proposition~\ref{prop:antecedent-state}) --- in particular, if a graph can mint its first distinctions without any seeded from its container.
\item \textsc{[shape]}~\textbf{Residual is inadmissibility.} Fails if a graph-originated deformation can be left as residual by a containing frame that nonetheless admits it under $\Gamma_{W_c}$, or absorbed while inadmissible, so that residual and non-admission come apart.
\item \textsc{[shape]}~\textbf{Realisation is single-valued.} Fails if realised dynamics are better accounted for by retaining the admissible set than by selecting one member of it --- and the qualification is the whole test: a treatment carrying all continuations forward must recover the observed statistics from the retained structure \emph{without adding a measure over it to do so}, since retention that needs a weighting has reintroduced selection under another name.
\item \textsc{[shape]}~\textbf{No-copy} (the destroy\slash preserve postulate's exposure)\textbf{.} Fails if a substrate change can preserve its predecessor
without paying to keep it --- if retention is available at no cost rather than as a purchased record.
The retention-cost postulate requires a specified resource ledger. Landauer bounds logically irreversible erasure and does not require every change to be irreversible; reversible uncomputation is compatible with it (Section~7 of the long paper (\emph{The frame interface and its economics})). The arrow of the event order is \emph{not} at stake, though it is
easily thought to be: direction is carried by the asymmetry of enablement, which stands
whether or not a change preserves what it followed.
\item \textsc{[shape]}~\textbf{Decoherence is one operation.} Fails if the manner in which an inadmissible write fails is a property of what was attempted rather than of the frame's rules --- that is, if one frame under one rule set yields structurally distinct failure modes for two inadmissible writes, with the difference not attributable to that frame's own policy (Proposition~\ref{prop:decoherence-one}). Different outcomes across \emph{different} frames do not bear on this; that is what the claim predicts.
\item \textsc{[shape]}~\textbf{A swallowed failure is relocated, never closed.} Fails if a frame can \emph{settle} from inside that a swallowed failure went nowhere, rather than relocating the question to what it cannot inspect (Proposition~\ref{prop:handler-relocates}). Not that a frame cannot detect its own handler discarding a failure --- it can; the exposure is the settling. One of the four places named below at which the same event would be found.
\item \textsc{[shape]}~\textbf{Dependents are not held by the referent.} Fails if realised graph dynamics require a structure to carry the set of things depending on it --- back-references primitive rather than installed --- so that the consequences of removing a structure are recoverable from that structure alone, without querying the frames holding the references and without paid machinery (Proposition~\ref{prop:vacated-three-states}, argued at Remark~9.14 of the long paper (\emph{What never cohered is absent; what cohered and stopped leaves a hole})). A substrate in which deletion is always clean falsifies this, and would also remove the distinction between an absence and a vacancy.
\item \textsc{[shape]}~\textbf{Termination follows supply, not containment.} Fails if containment alone determines it --- if a subgraph must stop when the graph containing it stops, even where its primitives are supplied from beyond that graph (Proposition~\ref{prop:termination-supply}). The partial case carries its own clause: fails if a frame losing part of its supply must either stop or continue correctly, rather than running with an incomplete operation set and undefined behaviour, or if such a frame can certify exhaustive unrestricted knowledge of its lost supply from inside, without assuming a complete provider registry. A local enumeration under a declared registry does not falsify the provenance commitment below.
\item \textsc{[shape]}~\textbf{A boundary is known only by crossing it.} Fails if a frame can \emph{settle} the location of its own admissibility boundary without any write that crosses it. Extrapolation from inside the admissible region evidences a boundary and is not at issue; establishing one would be (Proposition~\ref{prop:boundary-by-crossing}). This is the epistemic member of the same family as the provenance commitment below, and fails in the same manner --- by a procedure that closes the question rather than relocating it.
\item \textsc{[detail]}~\textbf{Robustness trades against legibility.} Fails if a policy can bound a frame's residuals without reducing what its failures disclose about its grammar --- a handler that both absorbs and reports with no loss (Proposition~\ref{prop:boundary-by-crossing}, argued at Remark~9.12 of the long paper (\emph{Forcing a failure is the instrument, and handlers blunt it})). Cheaply testable in software, and marked \textsc{[detail]} because the account survives its loss with the trade weakened to a tendency.
\item \textsc{[shape]}~\textbf{Capability and accessibility are distinct.} Fails if what a frame's gates admit and what it can reach from where it stands collapse into one quantity --- if there is no case of a frame able to represent and perform an arrangement it cannot reach, or none of a frame able to reach what it cannot represent (Definition~\ref{def:capability}). The coupling claimed between them, that a capacity shortfall narrows aperture and so degrades capability over time, fails separately if aperture width is unaffected by demands on capacity.
\item \textsc{[detail]}~\textbf{A hazard is a signed capability change above a frame-relative threshold.} Fails if a disclosure can be a hazard to a frame in which the negative sign is not admissible, or if the magnitude at which one counts is fixed independently of the frame bearing it rather than as a ratio to what that frame can absorb. Fails also if the threshold sits somewhere other than where the impact forces a write.
\item \textsc{[shape]}~\textbf{The relativity of grounding is a consequence, and can fail as one.} Fails if a view can be informative without standing in relation to what it views, or if relation is not what carries information here. A consequence has content and can fail, which a posit or a costless guard could not (Remark~2.5 of the long paper (\emph{The axiom in bits})). The second disjunct is \emph{Is}'s own identification of distinction with relation, so this is where the account is thinnest rather than where it is safest, and it is marked \textsc{[shape]} because what fails with it is the range of everything derived through it.
\end{enumerate}

\medskip\noindent\emph{Not carried in this statement:} A held program is one write; Coupling is a setting, not a kind; Globally-async/locally-sync; Phase is signed; substrate time is not; Externalising lengthens the interval before a crossing is legible; No free lunch; Description length is constitutive, not merely a method.



\noindent\small\emph{The works named above, listed from the long paper so the two cannot disagree.}

\begin{thebibliography}{99}
\itemsep=2pt

\bibitem[Buneman et~al.(2001)]{buneman2001}
Peter Buneman, Sanjeev Khanna, and Wang-Chiew Tan.
\newblock Why and where: A characterization of data provenance.
\newblock In Jan Van den Bussche and Victor Vianu, editors,
\emph{Database Theory --- ICDT 2001}, Lecture Notes in Computer Science
1973, pages 316--330. Springer, 2001.
\newblock \href{https://doi.org/10.1007/3-540-44503-X_20}{doi:10.1007/3-540-44503-X\_20}.

\bibitem[Johnson et~al.(1993)]{johnson1993}
Marcia K. Johnson, Shahin Hashtroudi, and D. Stephen Lindsay.
\newblock Source monitoring.
\newblock \emph{Psychological Bulletin}, 114(1):3--28, 1993.
\newblock \href{https://doi.org/10.1037/0033-2909.114.1.3}{doi:10.1037/0033-2909.114.1.3}.

\bibitem[Shpitser and Pearl(2008)]{shpitserpearl2008}
Ilya Shpitser and Judea Pearl.
\newblock Complete identification methods for the causal hierarchy.
\newblock \emph{Journal of Machine Learning Research}, 9:1941--1979, 2008.
\newblock \url{https://jmlr.org/papers/v9/shpitser08a.html}.

\end{thebibliography}
\end{document}


